\chapter{Managed and Functional Languages on the Desk}\label{ch:ll:managed-and-functional-languages-on-the-desk}

In 2011 a retail exchange described in public how its matching logic processed six million orders a second on a single
thread, written in Java (Fowler, 2011). Its engineers had not escaped the \omterm{def:ll:managed-and-functional-languages-on-the-desk:gc}{garbage collector}; they had designed around it:
everything pre-allocated, ring buffers of millions of slots created at start-up, no locks in the business logic. A few
years earlier a proprietary trading firm had described why it wrote its critical trading systems in OCaml, a functional
language with a \omterm{def:ll:managed-and-functional-languages-on-the-desk:gc}{garbage collector} of its own (Minsky and Weeks, 2008). Managed languages are on trading desks, sometimes on
the hot path. This chapter explains what their runtimes do that C++ and Rust do not, how a latency-sensitive program
written in them is structured, and what the firms that use them have said in public about why.

\section{Garbage collection and its pauses}

\begin{definition}[Garbage collector, stop-the-world pause]\label{def:ll:managed-and-functional-languages-on-the-desk:gc}
A \emph{garbage collector}\index{garbage collector} reclaims memory that the program can no longer reach, by tracing
references from the program's roots (stacks, globals). Most divide the heap into generations: new objects in a small young
generation, collected often and cheaply; survivors promoted to an old generation, collected rarely and at a cost that grows
with its size. A \emph{stop-the-world pause}\index{stop-the-world pause} is an interval during which the collector has
stopped every application thread.
\end{definition}

The trade is explicit: the program never frees memory and never frees it twice; in exchange, it runs on a schedule set by
its allocation. Modern collectors shorten pauses by doing most of their work concurrently with the program; a widely used
Java collector aimed first at pauses under ten milliseconds and later at under one millisecond spent at the points where
all threads stop (dated box below). Either figure is enormous against a microsecond budget, and the concurrent work itself
takes cores and memory bandwidth from the program.

\begin{omfigure}
\begin{tikzpicture}[font=\footnotesize]
\draw[fill=omDef!10,draw=omDef,rounded corners=2pt] (0,0) rectangle (3.2,1.2);
\node[anchor=north] at (1.6,1.15) {young generation};
\node at (1.6,0.45) {new objects};
\draw[fill=omThm!10,draw=omThm,rounded corners=2pt] (5.0,0) rectangle (11.2,1.2);
\node[anchor=north] at (8.1,1.15) {old generation};
\node at (8.1,0.45) {long-lived objects: the book, caches, state};
\draw[-Stealth,thick] (3.3,0.6) -- node[above]{survivors} node[below]{promoted} (4.9,0.6);
\draw[-Stealth,thick] (-1.4,0.6) node[left]{allocation} -- (-0.1,0.6);
\draw[omDef,thick,decorate,decoration={brace,amplitude=4pt,mirror}] (0,-0.15) -- (3.2,-0.15) node[midway,below=5pt,align=center]{collected often:\\cost $\propto$ survivors};
\draw[omThm,thick,decorate,decoration={brace,amplitude=4pt,mirror}] (5.0,-0.15) -- (11.2,-0.15) node[midway,below=5pt,align=center]{collected rarely:\\cost grows with the live heap};
\end{tikzpicture}

\omcaption{A generational heap. Allocation fills the young generation, whose collections cost in proportion to what
survives; survivors move to the old generation, whose full collections scan everything alive, and stop or slow the
program for longer as the heap grows.}\label{fig:ll:managed-and-functional-languages-on-the-desk:heap}
\end{omfigure}

\begin{dated}{2026-09}{Pause goals of a modern Java collector}\label{dat:ll:managed-and-functional-languages-on-the-desk:jvm}
OpenJDK's ZGC was proposed (JEP 333) with the goal that garbage-collection pause times should not exceed 10~milliseconds,
its stop-the-world phases limited to scanning roots. A later proposal (JEP 376) moved the scanning of thread stacks to a
concurrent phase, with the expectation that less than one millisecond be spent inside the collector's safepoints on typical
machines.
\end{dated}

\section{The low-latency managed style: pre-allocation, off-heap memory, warm-up}

\begin{definition}[Off-heap memory]\label{def:ll:managed-and-functional-languages-on-the-desk:offheap}
\emph{Off-heap memory}\index{off-heap memory} is memory a managed program uses outside its collected heap (direct byte
buffers, memory-mapped files, native arrays), laid out by the program and invisible to the collector: it is never scanned,
moved or freed by it.
\end{definition}

The style that the exchange above described follows from the collector's economics. Allocate nothing per message after
start-up: pre-allocate every object the hot path needs and reuse it (the pools of chapter~6); keep large, long-lived data
(order books, rings, journals) off the heap or in flat arrays of primitives, so that full collections have little to scan;
size the young generation so that the little garbage that remains fills it rarely; and run a quiet collector. What is
left is a program whose collector has nothing to do during the trading day. The limit of the style is its brittleness:
one library call that allocates, one boxed integer, puts the collector back on the path, and nothing in the language says
so; the allocation must be measured, as in chapter~6.

\begin{proposition}[The garbage budget]\label{prop:ll:managed-and-functional-languages-on-the-desk:budget}
A program that allocates $b$ bytes per message at $r$ messages a second, with a young generation of $Y$ bytes, collects its
young generation every $Y/(rb)$ seconds. It collects at most $k$ times in a session of length $T$ if and only if
$b \le kY/(rT)$.
\end{proposition}

\begin{proof}
Allocation fills the young generation at $rb$ bytes a second and each collection empties it, so collections occur every
$Y/(rb)$ seconds and number $rbT/Y$ in the session; $rbT/Y \le k$ is the stated condition.
\end{proof}

With a young generation of 2~GiB, 200\,000 messages a second and a session of six and a half hours, the budget for one
collection a session is under half a byte per message: in practice, zero.

\section{Just-in-time compilation and deoptimisation}

\begin{definition}[Just-in-time compilation, warm-up period]\label{def:ll:managed-and-functional-languages-on-the-desk:jit}
In \emph{just-in-time compilation}\index{just-in-time compilation} a runtime first interprets or lightly compiles a program,
measures which code is hot and how it behaves, and then compiles the hot code to optimised machine code using that profile.
The \emph{warm-up period}\index{warm-up period} is the time, or number of executions, before the hot path runs as optimised
code.
\end{definition}

\begin{definition}[Deoptimisation]\label{def:ll:managed-and-functional-languages-on-the-desk:deopt}
\emph{Deoptimisation}\index{deoptimisation} is the runtime's fall-back from optimised code to a slower tier when an
assumption the optimiser made from the profile is violated (a call site that had seen one type sees another, a branch never
taken is taken); the code may later be recompiled.
\end{definition}

A just-in-time compiler is \omterm{def:ll:cpp-for-latency-iii-the-compiler:pgo}{profile-guided optimisation} (chapter~8) done while the program runs, with the same risk: the
profile is whatever the program has seen so far. A trading process that starts in the morning is slowest on its first
messages, exactly when the open is busiest, and a rare path taken for the first time at a volatile moment (an unusual
message type, a halt) may trigger \omterm{def:ll:managed-and-functional-languages-on-the-desk:deopt}{deoptimisation} just then. The standard remedies are warm-up by replaying synthetic or
recorded traffic through the real code before trading (with orders sent to a simulator, never the market), and keeping
the hot path's types and branches stable. Ahead-of-time compilation removes the warm-up at the cost of the profile.

\section{The functional-language shops}

The proprietary firm of the opening described its choice in a journal paper: OCaml as its primary development language,
over twenty OCaml programmers and hundreds of thousands of lines of code, used for critical trading systems, research,
systems software and administration, valued because it let the firm ``rapidly produce readable, correct, efficient code''
(Minsky and Weeks, 2008). The argument was not speed per message but correctness per engineer: an expressive type system
that encodes invariants (an order's state as a variant whose cases the compiler checks are all handled), immutable data by
default, and code short enough to be reviewed by people who trade with it. The same firm's later essay argued that the
next language a programmer learns should be functional (Minsky, 2011). OCaml's collector is generational too, with a small
minor heap collected often and a major heap collected incrementally; latency-critical OCaml, like latency-critical
Java, is written to allocate little on the hot path.

\section{Choosing a language for a desk}

The choice is less binary than debates suggest. Several languages usually coexist in one firm: C++ or Rust where
nanoseconds are traded, a managed or functional language for the large body of logic whose
\omterm{def:ll:where-latency-comes-from:budget}{latency budget} is microseconds to milliseconds, and Python for research and operations (One Quant Book~15, chapter~8). The
questions that decide are measurable: the \omterm{def:ll:where-latency-comes-from:budget}{latency budget} at the percentile that matters; the allocation discipline the
team can sustain and test; the cost of the bugs each language prevents or permits; and the size and experience of the
team. In Python, this book's research language, the lessons of this chapter apply directly, as the tutorial shows.

\section{Tutorial: the collector in a Python message loop}\label{tut:ll:managed-and-functional-languages-on-the-desk:tut}

\begin{tutorial}
\textbf{Goal.} Watch CPython's collector interrupt a message loop, and remove it by pre-allocation. \textbf{End state:}
\cref{fig:ll:managed-and-functional-languages-on-the-desk:gc}.
\begin{tutsteps}
\item \textbf{Two books.} One keeps each live order as a Python dictionary, adding two orders for each one removed, as a
book fills during the morning; the other writes orders into a pre-allocated NumPy record array.
\omcode{code/low-latency/10-managed-and-functional-languages-on-the-desk/python/ll_gc.py}{19}{38}{A book of dictionaries and a pre-allocated book.}\label{lst:ll:managed-and-functional-languages-on-the-desk:books}
\item \textbf{Time the collector.} \texttt{GcWatch} registers a callback that CPython calls before and after every
collection, and records each pause, by generation, in \texttt{firm.lathist}.
\omcode{code/firm/gcwatch/firm_gcwatch.py}{27}{50}{Timing every collection through the gc module's callbacks.}\label{lst:ll:managed-and-functional-languages-on-the-desk:watch}
\item \textbf{Measure the loop} over 200\,000 messages with \texttt{python bench\_gc.py}; the loop stores timestamps only,
so that nothing but the book's own work happens between two of them.
\end{tutsteps}
\textbf{What to change next.} Keep the dictionary book but call \texttt{gc.freeze()} after the warm-up and again every
100\,000 messages, and see what happens to the full collections; bound the dictionary book at 1\,000 live orders and
explain why the collector then never runs.
\end{tutorial}

\begin{omfigure}
\begin{tikzpicture}
\begin{axis}[width=10cm,height=5.6cm,ybar,bar width=8pt,ymode=log,log origin y=infty,ymin=100,ymax=100000000,
  symbolic x coords={dict book,pooled book},xtick=data,enlarge x limits=0.45,
  ylabel={ns per message (log scale)},tick label style={font=\small},label style={font=\small},grid=major,
  grid style={gray!25},legend columns=4,legend style={font=\footnotesize,at={(0.5,-0.2)},anchor=north,draw=none,
  fill=none,/tikz/every even column/.append style={column sep=6pt}},table/col sep=comma]
\addplot[fill=omDef!70,draw=omDef] table[x=design,y=p50]{figdata/low-latency/10-managed-and-functional-languages-on-the-desk/measured_gc.csv};
\addplot[fill=omDef!35,draw=omDef] table[x=design,y=p99]{figdata/low-latency/10-managed-and-functional-languages-on-the-desk/measured_gc.csv};
\addplot[fill=omThm!45,draw=omThm] table[x=design,y=p999]{figdata/low-latency/10-managed-and-functional-languages-on-the-desk/measured_gc.csv};
\addplot[fill=gray!40,draw=gray] table[x=design,y=max]{figdata/low-latency/10-managed-and-functional-languages-on-the-desk/measured_gc.csv};
\legend{p50,p99,p99.9,maximum}
\end{axis}
\end{tikzpicture}

\omcaption{A Python message loop over 200\,000 orders: a book of dictionaries, which triggered 259 young, 23 intermediate
and 2 full collections, and a pre-allocated NumPy book, which triggered none. The dictionary book's maximum is its full
collection. CPython 3.10. Measured on a laptop (Intel Core Ultra 7 155H) under WSL2, no isolated cores. Data:
\texttt{bench\_gc.py}.}\label{fig:ll:managed-and-functional-languages-on-the-desk:gc}
\end{omfigure}

The result has two halves, and both matter. The pre-allocated book removed the collector: no collection in 200\,000
messages, and a tail more than twenty times shorter at its maximum. It is not faster at the median: writing three fields
of a NumPy array from Python costs more than creating a small dictionary. And the dictionary book's worst message took as
long as its last full collection, some twenty milliseconds, which scanned every order still alive. CPython adds a
twist of its own: its collector counts allocations minus deallocations, so a loop whose live set stays bounded never
triggers a collection at all; a book that grows does.

\section{Build: collector watch and message pool}\label{bld:ll:managed-and-functional-languages-on-the-desk:gcwatch}

\begin{build}
\bfield{Purpose} The firm's Python services (research servers, reporting, the simulator drivers of One Quant Book~10)
measure their collector and keep it off their message loops.
\bfield{Interface} Python \texttt{firm\_gcwatch}: \texttt{GcWatch()} as a context manager with \texttt{pauses[gen]}
(\texttt{LatHist}, nanoseconds), \texttt{collections[gen]}, \texttt{total\_ns}; \texttt{RecordPool(n, dtype)} with
\texttt{next()} and \texttt{buf}; \texttt{MESSAGE} (timestamp, price, quantity, side); \texttt{steady\_state(step, n,
warm)}.
\bfield{Rules} The callback is removed when the context ends, even on error; \texttt{steady\_state} freezes the heap after
warm-up and unfreezes it afterwards; the pool never grows.
\bfield{Acceptance tests} \texttt{code/firm/gcwatch/tests/}: a pooled loop of 50\,000 messages causes no collection; a
bounded churn causes none either; a growing book causes young collections, each timed.
\bfield{Stretch} Report the allocation rate (bytes a message) with \texttt{tracemalloc} in a test build; alarm when a
full collection happens during market hours.
\end{build}

\begin{omsources}
\item M.~Fowler, ``The LMAX architecture'', 12 July 2011.
\item Y.~Minsky and S.~Weeks, ``Caml trading: experiences with functional programming on Wall Street'', \emph{Journal of
Functional Programming} 18(4), 2008; Y.~Minsky, ``OCaml for the masses'', \emph{ACM Queue}, 2011.
\item OpenJDK, JEP 333 (ZGC) and JEP 376 (concurrent thread-stack processing).
\item Python documentation, module \texttt{gc}.
\end{omsources}

\section{Exercises}

\begin{exercise}[$\star$]\label{exo:ll:managed-and-functional-languages-on-the-desk:1}
A Java engine allocates 64 bytes per message at 100\,000 messages a second with a 1~GiB young generation. How often does it
collect the young generation, and how many times in a 6.5-hour session?
\end{exercise}

\begin{exercise}[$\star$]\label{exo:ll:managed-and-functional-languages-on-the-desk:2}
A collector's pause is \qty{1}{\milli\second}. How many messages arriving at 200\,000 a second queue behind it, and how long
does the last of them wait?
\end{exercise}

\begin{exercise}[$\star$]\label{exo:ll:managed-and-functional-languages-on-the-desk:3}
From \cref{fig:ll:managed-and-functional-languages-on-the-desk:gc}, by how much does the pooled book shorten the maximum and
the p99.9, and how does its median compare?
\end{exercise}

\begin{exercise}[$\star\star$]\label{exo:ll:managed-and-functional-languages-on-the-desk:4}
Why does a Python loop that keeps the last 1\,000 orders in a bounded deque never trigger the collector, while a growing book
does?
\end{exercise}

\begin{exercise}[$\star\star$]\label{exo:ll:managed-and-functional-languages-on-the-desk:5}
A Java strategy is fastest after ten minutes of trading and slowest at the open. Explain, and propose two remedies.
\end{exercise}

\begin{exercise}[$\star\star$]\label{exo:ll:managed-and-functional-languages-on-the-desk:6}
The exchange of the opening kept ring buffers of 20 million slots on the input side. At 6 million orders a second, how many
seconds of input does such a ring hold?
\end{exercise}

\begin{exercise}[$\star\star\star$]\label{exo:ll:managed-and-functional-languages-on-the-desk:7}
\emph{Coding.} Call \texttt{gc.freeze()} in the dictionary book after the warm-up and then every 100\,000 messages. Count the
full collections and record the worst message. What does freezing trade away?
\end{exercise}

\begin{exercise}[$\star\star\star$]\label{exo:ll:managed-and-functional-languages-on-the-desk:8}
\emph{Find the flaw.} ``Our Java gateway uses the concurrent collector with pauses under a millisecond, so garbage collection
does not affect our latency.''
\end{exercise}

\section{Problem: The Garbage Budget}

\begin{problem}[{Weekend problem --- a JVM engine on a busy day}]\label{pb:ll:managed-and-functional-languages-on-the-desk:1}
A firm runs an order-routing engine on the JVM with a 2~GiB young generation. On a normal day it handles 200\,000 messages a
second for a 6.5-hour session and allocates 48 bytes per message on average. A code review finds that most of it comes from
boxed integers and a logging call that formats a string.

\medskip\noindent\textbf{Part I --- The budget.}
\begin{enumerate}
\item How many bytes does the engine allocate per second, and per session?
\item How often does it collect the young generation, and how many times per session?
\item What is the allocation per message that allows at most one young collection per session?
\item What does that budget mean in practice?
\end{enumerate}

\medskip\noindent\textbf{Part II --- The pauses.}
\begin{enumerate}[resume]
\item If each young collection pauses the engine for \qty{0.5}{\milli\second}, how much time a session does it spend paused?
\item How many messages queue behind one such pause, and what does the one at the back of the queue wait?
\item On a day with three times the traffic, how do the answers change?
\item Why is the busy day the one that matters?
\end{enumerate}

\medskip\noindent\textbf{Part III --- The fixes.}
\begin{enumerate}[resume]
\item What replaces boxed integers on the hot path?
\item What replaces formatted logging (chapter~23)?
\item Where should the order book and the rings live, and why?
\item How would you verify, in the test suite, that the allocation per message is zero?
\end{enumerate}

\medskip\noindent\textbf{Part IV --- The verdict.}
\begin{enumerate}[resume]
\item State the \emph{named result}: the allocation per message that keeps the engine to one young collection a session,
and the pause time a session at the current 48 bytes.
\item What remains of the collector's cost once allocation is zero?
\item Why must the engine still be warmed up before the open?
\item What did the Python experiment of the tutorial show that carries over to the JVM?
\item Why did the functional-language firm value its language, according to its own paper?
\item When would you choose C++ or Rust over a managed language for this engine?
\item What must the firm measure before deciding?
\item In one sentence: what does a \omterm{def:ll:managed-and-functional-languages-on-the-desk:gc}{garbage collector} cost a hot path?
\end{enumerate}
\end{problem}

\section{Interview questions}

\begin{interviewq}[$\star$ \iqroles{developer}]\label{iq:ll:managed-and-functional-languages-on-the-desk:1}
How does a generational \omterm{def:ll:managed-and-functional-languages-on-the-desk:gc}{garbage collector} work, and why is it generational?
\end{interviewq}

\begin{interviewq}[$\star\star$ \iqroles{developer}]\label{iq:ll:managed-and-functional-languages-on-the-desk:2}
How would you write a low-latency service in Java so that the collector does not interfere?
\end{interviewq}

\begin{interviewq}[$\star\star$ \iqroles{developer}]\label{iq:ll:managed-and-functional-languages-on-the-desk:3}
What is warm-up in a JIT-compiled runtime, and how do you deal with it before the open?
\end{interviewq}

\begin{interviewq}[$\star\star$ \iqroles{developer}]\label{iq:ll:managed-and-functional-languages-on-the-desk:4}
What is \omterm{def:ll:managed-and-functional-languages-on-the-desk:deopt}{deoptimisation}, and why can it happen at the worst moment?
\end{interviewq}

\begin{interviewq}[$\star\star$ \iqroles{developer, researcher}]\label{iq:ll:managed-and-functional-languages-on-the-desk:5}
Why might a trading firm choose a functional language for its trading systems?
\end{interviewq}

\begin{interviewq}[$\star\star\star$ \iqroles{developer}]\label{iq:ll:managed-and-functional-languages-on-the-desk:6}
A Python service's p99.9 is fine but its maximum is 20 milliseconds once every few minutes. How do you find out whether it is
the \omterm{def:ll:managed-and-functional-languages-on-the-desk:gc}{garbage collector}, and what do you do about it?
\end{interviewq}
